\documentclass[11pt,a4paper]{article}
\usepackage[margin=23mm]{geometry}
\usepackage[T1]{fontenc}
\usepackage[utf8]{inputenc}
\usepackage{lmodern}
\usepackage{microtype}
\usepackage{longtable,booktabs,array}
\usepackage{hyperref}
\usepackage{enumitem}
\setlist[enumerate]{leftmargin=*}
\setlength{\parskip}{0.5em}
\setlength{\parindent}{0pt}
\title{Assessment of the IEEE TAC Journal Reviews}
\date{7 October 2026}
\begin{document}
\maketitle
This assessment compares the associate editor's report and the comments of Reviewers 3, 4, and 5 with the \textbf{current} \texttt{main.tex} (7 October 2026). The submitted manuscript is present as \texttt{26-0591\_01\_MS-1.pdf} (the file name differs from \texttt{26-0591\_01\_MS-21.pdf} in the request). The current manuscript contains intervening changes described in \texttt{response\_to\_reviewers.tex}; that response addresses the earlier mathematical review, \textbf{not} the journal reviews. I use it only as context for changes already made. Status means whether the *journal comment* is answered in the current manuscript, not whether the mathematics has been independently validated.

\section{Overall assessment}
The present revision has improved several assumptions and corrected some examples, but it is \textbf{not yet ready for resubmission} against the journal reviews. The editor's central concerns remain: the original contribution is insufficiently separated from prior results, Theorems 2–3 lack self-contained proofs or exact proof references, and the model and algebraic algorithms need much clearer explanation. Reviewer 5's concerns about length and focus also remain. The editor suggests rejection with an option to resubmit after major revision, and warns that the next round must be close enough to acceptance to require at most minor changes. This makes substantive revision and a point-by-point journal response important.

\section{Comment-by-comment assessment}
\small
\begin{longtable}{@{}p{0.22\textwidth}p{0.32\textwidth}p{0.40\textwidth}@{}}
\toprule
\textbf{Journal comment} & \textbf{Status in current manuscript} & \textbf{Revision needed}\\
\midrule
\endfirsthead
\toprule
\textbf{Journal comment} & \textbf{Status in current manuscript} & \textbf{Revision needed}\\
\midrule
\endhead
\textbf{R3.1: Conditions under which `P` is invertible} & \textbf{Partial.} Section II, lines 132–145, now explicitly assumes \texttt{P(q)} is invertible for almost all \texttt{q}; lines 372–376 explain a row permutation for generic invertibility. Neither states usable structural or analytic conditions for well-posedness or gives a reference. & Define precisely whether invertibility means over rational functions, pointwise on the unit circle, or both; distinguish it from stability. Give a familiar sufficient condition for \texttt{P=I-G} and explain how the assumption is checked in the proposed general model.\\\addlinespace\midrule
\textbf{R3.2 / R4: formal proofs of informativity theorems} & \textbf{Partial.} Theorem 1 has a short proof at lines 316–337 and now requires positive definite spectra for almost all frequencies. Theorem 2 at lines 434–439 has no formal proof; lines 385–408 contain a proof sketch and an appeal to \texttt{wou91}. Theorem 3 at lines 473–489 has no proof. & Give full proofs for Theorems 2–3, including the direct-feedthrough extension, genericity and free-entry assumptions, and the saturation/elimination cases of Theorem 3; alternatively identify an exact existing theorem and show explicitly that every hypothesis maps to this model. Expand Theorem 1's proof where the generic-rank-to-almost-all-frequency step is used.\\\addlinespace\midrule
\textbf{R3.3: Which predictors can be selected?} & \textbf{Partial.} Section III constructs a predictor; lines 348–359 clarify when predictor inputs can be reduced, and Example 1 gives a choice. There is no general selection procedure or statement of how a chosen predictor changes the identifiability target. & State the class of allowed predictors, how measured outputs and predictor inputs are chosen, and which transfer functions remain uniquely recoverable. Give a small selection example or practical decision rule.\\\addlinespace\midrule
\textbf{R3.4: Is topology known?} & \textbf{Open.} The graph uses structural zeros (lines 410–449), but the assumptions do not explicitly say which sparsity pattern is known and which edges or coefficients are unknown. & State the prior topology/sparsity assumptions separately for the graphical free-entry result and the constrained Gröbner basis result; clarify whether topology recovery is outside scope.\\\addlinespace\midrule
\textbf{R3.5–6: How to construct the Gröbner bases for Theorems 3–4?} & \textbf{Partial.} Lines 456–470 define minors, denominators and the ideal for Theorem 3; lines 634–648 sketch elimination for Theorem 4. Variable orders are mentioned. The paper still lacks an executable small example, explicit variable lists and a clear interpretation of each output. & Add a compact worked calculation or algorithm box for each test: parameterize free entries, build denominator exclusions and ideals, specify monomial order, compute the reduced basis/elimination ideal, then interpret rank or image dimension. Clarify coefficient field and how equality constraints and reducible domains are handled.\\\addlinespace\midrule
\textbf{R4: Distinguish ordinary, network, generic, local and global identifiability} & \textbf{Partial.} Lines 572–599 now define generic local *algebraic* network identifiability and distinguish finite ambiguity from generic uniqueness. But the introduction, lines 101–104, still calls uniqueness of open-loop functions from a predictor simply “identifiability,” and the relationship to classical parameter identifiability remains unclear. & Define the terms near the start; consistently use “network identifiability” for the map from network parameters to observed transfers; distinguish generic uniqueness, local finite-fiber behavior and pointwise uniqueness. Explain the distinction between symbolic rational-function and single-frequency definitions.\\\addlinespace\midrule
\textbf{R4: Partial measurement is not new in itself} & \textbf{Mostly addressed.} The introduction now cites prior subnetwork work at lines 93–95 and claims a particular generalization at lines 107–128, rather than saying all previous work used full measurement. & Tighten “more general than previously possible” to precisely named combinations of assumptions, with a comparison to closest prior work. Avoid implying that partial measurement alone is the novelty.\\\addlinespace\midrule
\textbf{R4 / R5: Informativity versus identifiability; informativity is standard MIMO} & \textbf{Partial.} Sections IV and V are separate and the introduction describes two stages. The informativity section still presents much of the rank argument as a network contribution, without clearly saying that the excitation criterion is a standard predictor/MIMO criterion applied to a network-derived transfer matrix. & State the logical pipeline and independent scope of each result. Attribute standard informativity and transfer-of-excitation facts; identify exactly which part depends on network structure. Consider shortening the generic informativity material if novelty is limited.\\\addlinespace\midrule
\textbf{R4: Reference for Definition 1 and other inherited results} & \textbf{Open/partial.} The definition at lines 242–253 has no citation. Some inherited graph and dimension results have citations, but several claims are described as generalizations without exact source statements. & Cite the precise source for the informativity definition (or explain any difference); audit definitions and propositions for attribution. Include a theorem-by-theorem comparison table or short paragraph.\\\addlinespace\midrule
\textbf{R4: “Almost all frequencies” in Theorem 1} & \textbf{Addressed.} Lines 318–320 and 327–335 state and use this qualification. & Preserve this wording and make the measure/genericity conventions precise.\\\addlinespace\midrule
\textbf{R4: Gevers–Bazanella–Bombois–Miskovic transfer-of-excitation context} & \textbf{Open.} The requested 2009 IEEE TAC paper exists in \texttt{refs.bib} as key \texttt{5325719}, but \texttt{main.tex} never cites it. & Cite it in the informativity discussion and compare its assumptions and claim to Theorem 1.\\\addlinespace\midrule
\textbf{R4: Novelty and proof of Theorem 2; proof/newness of Theorem 3} & \textbf{Open/partial.} Theorem 2 is explicitly called a generalization of \texttt{wou91} (lines 431–432), with only a brief extension argument. Theorem 3 is stated without proof or a precise novelty claim. & State which results are known, which hypotheses are widened (especially direct feedthrough, partial measurement and constraints), and what technical step makes each extension valid. If a theorem is a restatement, label and cite it as such.\\\addlinespace\midrule
\textbf{R4: `z` used twice} & \textbf{Open.} \texttt{z\_k} is an output (line 137), while \texttt{z} is also the complex transform variable (line 146 and throughout the algebraic sections). & Rename one of these consistently, including equations and examples.\\\addlinespace\midrule
\textbf{R5.1: Why the `P,Q,R,S` model rather than standard network models?} & \textbf{Open.} Lines 132–145 give dimensions, and later sections use the model, but there is no motivating derivation or explicit embedding of \texttt{w=Gw+Br+v}, \texttt{y=Cw}. & Show the standard model as \texttt{P=I-G} with corresponding \texttt{Q,R,S}, explain the roles of each matrix, and give one concrete case requiring dynamic \texttt{Q} or \texttt{R} and direct \texttt{S}. Explain the state/output split and any restrictions needed for physical network interpretation.\\\addlinespace\midrule
\textbf{R5.2–3 / editor: results look straightforward; path/rank link already known} & \textbf{Open/partial.} The current text cites \texttt{wou91} and claims extensions, but still does not demonstrate why the contribution is nontrivial or compare directly to manuscript reference [3] as the reviewer requested. & Map \textbf{submitted-paper reference numbers} [3], [17] and other nearest results to current bibliography entries before drafting the response. Compare assumptions and conclusions side by side. Reframe known path/rank equivalence as a tool unless a new theorem with a proved extension can be justified.\\\addlinespace\midrule
\textbf{R5.5–6: Remove/condense examples and subnetwork section} & \textbf{Open.} Sub-Networks (lines 699–957), Example (958–1115) and Numerical Examples (1117–1190) remain extensive. The “11 minutes on a laptop” sentence is still at line 1144. & Restructure around the strongest new result, fold the subnetwork case into the general setup where possible, keep one example that demonstrates a capability unavailable in closest prior work, and move or remove routine or weak numerical examples. If timings remain, report software, hardware, input sizes and reproducible code.\\\addlinespace\midrule
\textbf{R5.7: Example 3 versus cited conjecture; does two solutions mean unidentifiable?} & \textbf{Partial.} Lines 1073–1113 now distinguish generic local algebraic identifiability (two isolated solutions) from generic uniqueness and narrow the conjecture claim to equality-constrained duplicated networks. & Check the constructed second solution against every admissibility condition: nonzero edges, properness, stability, well-posedness and the exact assumptions of the cited conjecture. Prove generic degree two and show that the duplicated network is genuinely in the conjecture's model class. If it is not, present this as a limited counterexample rather than a contradiction.\\\addlinespace\midrule
\textbf{R5.8: `F(X)` does not display `X`} & \textbf{Mostly addressed.} Lines 531–545 now explicitly define \texttt{X} as the vector of free entries of \texttt{(P,Q\_r,R\_y,S\_yr)}. & For readability, write \texttt{P(X),Q\_r(X),R\_y(X),S\_\{yr\}(X)} in the definition of \texttt{F}, at least once.\\\addlinespace\midrule
\textbf{R5.9: Example 1 exceptional zero edge} & \textbf{Addressed in substance.} Lines 1025–1034 now say the result is *generic* network identifiability, with uniqueness failing at \texttt{G\textasciicircum{}1=0}. & State whether the zero transfer is admissible under the chosen fixed-topology convention. If edges are required to be structurally nonzero, avoid calling the zero case an admissible failure; if zero coefficients are allowed, explain the difference between structural edges and a nongeneric zero realization.\\\addlinespace\midrule
\textbf{R5.10: Introduce algebraic geometry terminology} & \textbf{Partial.} Sections IV–V use ideals, affine varieties, fibers, closures, dimensions and Gröbner bases, and a few are defined locally (lines 547–599). The paper still assumes substantial background. & Add a concise reader-oriented glossary and a one-variable or two-variable example illustrating ideal, variety, fiber, elimination and generic point; keep formal definitions sufficient to understand both algorithms.\\\addlinespace\midrule
\textbf{R5.11: Typos and bibliography} & \textbf{Partial.} The cited \texttt{2913} year appears corrected to \texttt{2013} in \texttt{refs.bib}, but the manuscript still has typographical/grammar issues (for example lines 521, 665, 1120–1122, 1145). & Perform a full language, notation, reference and typesetting pass after restructuring.\\\addlinespace\midrule
\bottomrule
\end{longtable}
\normalsize
\section{Recommended order of work}
\begin{enumerate}
\item \textbf{Set the claim and scope.} Decide which result is genuinely new and suitable for a technical note. Build a comparison of Theorems 1–4 with the exact cited prior theorems, including the sources denoted [3], [17] and [19] in the submitted version.
\item \textbf{Resolve the mathematical core.} Supply rigorous proofs or precise reductions for Theorems 2–3, strengthen the steps of Theorem 1, and verify Example 3 against the cited conjecture and admissibility assumptions. Revise or withdraw claims that do not survive this check.
\item \textbf{Specify the model and procedure.} Explain \texttt{P,Q,R,S}, topology and well-posedness assumptions, predictor selection, and worked Gröbner basis construction.
\item \textbf{Reorganize and shorten.} Separate the standard informativity background from the network-specific contribution, consolidate subnetwork treatment, and retain examples that establish the distinct value of the method.
\item \textbf{Finish the editorial pass.} Fix terminology and \texttt{z} notation, add the requested citations, update references, and draft a point-by-point response to \textbf{each journal reviewer and the editor} based on the revised text.
\end{enumerate}
The status judgments above are based on the source text. They do not certify the proofs, examples, numerical calculations or fit with IEEE TAC's length and resubmission rules.

\end{document}
